\documentclass[a4paper,12pt]{article}
\usepackage[T2A]{fontenc}
\usepackage[utf8]{inputenc}
\usepackage[russian]{babel}
\usepackage{amsmath,amssymb,mathtools}
\usepackage{icomma}
\usepackage[left=22mm,right=22mm,top=23mm,bottom=23mm]{geometry}
\usepackage{microtype}
\usepackage{tikz}
\usetikzlibrary{arrows.meta,positioning,shapes.geometric,calc}
\usepackage{tabularx,booktabs}
\usepackage{enumitem}
\usepackage{fancyhdr}
\usepackage{xcolor}
\usepackage{titlesec}
\usepackage{setspace}
\usepackage[bookmarks=false]{hyperref}
\definecolor{accent}{HTML}{008080}
\definecolor{darkaccent}{HTML}{00565A}
\definecolor{textgray}{HTML}{30383B}
\definecolor{muted}{HTML}{69777A}
\hypersetup{colorlinks=true,linkcolor=darkaccent,urlcolor=darkaccent,pdfauthor={Лёвин Артём Александрович},pdftitle={Чек-лист: отношения, пропорции и проценты}}
\color{textgray}
\setstretch{1.12}
\setlength{\parindent}{0pt}
\setlength{\parskip}{4pt}
\setlist[enumerate]{leftmargin=*,itemsep=3pt,topsep=4pt}
\setlist[itemize]{leftmargin=*,itemsep=2pt,topsep=3pt}
\titleformat{\section}{\normalfont\large\color{darkaccent}}{\thesection.}{0.65em}{}[\vspace{6pt}\titlerule]
\titlespacing*{\section}{0pt}{12pt}{8pt}
\titleformat{\subsection}{\normalfont\normalsize\bfseries\color{darkaccent}}{}{0pt}{}
\titlespacing*{\subsection}{0pt}{7pt}{3pt}
\pagestyle{fancy}
\fancyhf{}
\fancyhead[L]{\small Отношения, пропорции и проценты}
\fancyhead[R]{\small Ксения Васильченко}
\fancyfoot[C]{\small \thepage}
\renewcommand{\headrulewidth}{0.35pt}
\renewcommand{\headrule}{\hbox to\headwidth{\color{accent}\leaders\hrule height \headrulewidth\hfill}}
\setlength{\headheight}{26pt}
\newcommand{\key}[1]{\textcolor{darkaccent}{\textbf{#1}}}
\tikzset{
  flowstep/.style={draw=accent!85!black,rounded corners=3pt,line width=.8pt,align=center,text width=100mm,minimum height=15mm,inner sep=5pt,font=\normalsize},
  flowoval/.style={ellipse,draw=accent!85!black,line width=.8pt,align=center,text width=76mm,minimum height=12mm,inner sep=4pt,font=\normalsize},
  choice/.style={diamond,aspect=2.45,draw=accent!85!black,line width=.8pt,align=center,inner sep=3pt,font=\normalsize},
  answer/.style={draw=accent!85!black,rounded corners=3pt,align=center,text width=53mm,minimum height=18mm,inner sep=4pt,font=\normalsize},
  arr/.style={-{Stealth[length=3.4mm,width=2mm]},line width=.8pt,draw=darkaccent}
}
\begin{document}
\hypersetup{pageanchor=false}
\begin{titlepage}
\thispagestyle{empty}
\centering
\vspace*{27mm}
{\color{darkaccent}\Large МАТЕМАТИКА\par}
\vspace{21mm}
{\LARGE Чек-лист\par}
\vspace{15mm}
{\color{accent}\LARGE Деление числа в отношении,\\[3pt] пропорции и проценты\par}
\vspace{18mm}
{\normalsize Практические методы и самостоятельная тренировка\par}
\vfill
{\large Ксения Васильченко\par}
\vspace{12mm}
{\normalsize 19 ноября 2026 г.\par}
\vspace{18mm}
{\small Лёвин Артём Александрович\\ эксклюзивно для Ксении Васильченко\par}
\vspace{17mm}
\begin{tikzpicture}
\draw[accent,line width=.85pt] (0,0) .. controls (1.2,.52) and (2.4,-.30) .. (3.65,0)
  .. controls (4.95,.44) and (6.10,-.37) .. (7.55,0)
  .. controls (8.6,.27) and (9.2,-.26) .. (10,0);
\end{tikzpicture}
\vspace{8mm}
\end{titlepage}
\hypersetup{pageanchor=true}

\section{Как разделить число в заданном отношении}

Запись \(a:b\) показывает, \key{скольким долям соответствует каждая часть}. Если число \(N\) нужно разделить в отношении \(a:b\), где \(a,b>0\), действуйте так:
\begin{enumerate}
\item Найдите общее количество долей: \(S=a+b\).
\item Найдите величину одной доли: \(q=\dfrac{N}{S}\).
\item Найдите части: \(x=qa\) и \(y=qb\).
\item Проверьте, что \(x+y=N\).
\end{enumerate}

\key{Короткая запись:}
\[
\boxed{x=\frac{Na}{a+b},\qquad y=\frac{Nb}{a+b}.}
\]

\subsection{Пример 1. Отношение \(1:98\)}
Разделим число \(693\) в отношении \(1:98\).
\[
S=1+98=99,\qquad q=693:99=7.
\]
Следовательно,
\[
x=7\cdot1=7,\qquad y=7\cdot98=686.
\]
\key{Проверка:} \(7+686=693\); второе число в \(98\) раз больше первого.

\subsection{Пример 2. Отношение \(2:7\)}
Разделим число \(693\) в отношении \(2:7\):
\[
S=2+7=9,\qquad q=693:9=77;
\]
\[
x=77\cdot2=154,\qquad y=77\cdot7=539.
\]
Проверка: \(154+539=693\) и \(154:539=2:7\).

\section{Если частей больше двух}

Для отношения \(a:b:c\) идея остаётся той же: \(S=a+b+c\), а части равны \(qa\), \(qb\), \(qc\), где \(q=N/S\).

\subsection{Пример 3. Три части}
Разделим \(12\) в отношении \(1:2:3\):
\[
S=1+2+3=6,\quad q=12:6=2.
\]
\[
\boxed{2,\quad 4,\quad 6}\qquad (2+4+6=12).
\]
\clearpage
\thispagestyle{empty}
\begin{center}
{\Large\color{darkaccent} Алгоритм: разделение числа в отношении}
\end{center}
\vspace{14mm}
\begin{center}
\begin{tikzpicture}[node distance=10mm]
\node[flowoval] (start) {Дано число \(N\) и отношение \(a:b\) или \(a:b:c\)};
\node[flowstep,below=of start] (sum) {Сложите коэффициенты отношения\\ \(S=a+b\) или \(S=a+b+c\)};
\node[flowstep,below=of sum] (one) {Найдите величину одной доли\\ \(q=N:S\)};
\node[flowstep,below=of one] (parts) {Умножьте \(q\) на каждый коэффициент\\ \(qa,\ qb\) или \(qa,\ qb,\ qc\)};
\node[flowstep,below=of parts] (check) {Проверьте: сумма частей равна \(N\),\\ их отношение совпадает с условием};
\node[flowoval,below=of check] (done) {Ответ: найденные части};
\draw[arr] (start) -- (sum);
\draw[arr] (sum) -- (one);
\draw[arr] (one) -- (parts);
\draw[arr] (parts) -- (check);
\draw[arr] (check) -- (done);
\end{tikzpicture}
\end{center}
\clearpage

\section{Что такое пропорция}

\key{Пропорция} --- это верное равенство двух отношений. Её можно записать двумя способами:
\[
\boxed{a:b=c:d\qquad\Longleftrightarrow\qquad\frac{a}{b}=\frac{c}{d}},
\quad b\ne0,\ d\ne0.
\]
В пропорции \(a:b=c:d\) числа \(a\) и \(d\) называют \key{крайними членами}, а \(b\) и \(c\) --- \key{средними}.

\key{Основное свойство пропорции:}
\[
\boxed{\frac{a}{b}=\frac{c}{d}\quad\Longleftrightarrow\quad a\cdot d=b\cdot c}
\qquad (b,d\ne0).
\]
Это значит: произведение крайних членов равно произведению средних.

\subsection{Пример 4. Проверка верного равенства}
\[
\frac{2}{3}=\frac{4}{6}\ ?\quad
2\cdot6=12,\qquad 3\cdot4=12.
\]
Произведения совпали. \key{Вывод:} равенство является пропорцией.

\subsection{Пример 5. Опровержение}
\[
\frac{2}{3}=\frac{4}{5}\ ?\quad
2\cdot5=10,\qquad 3\cdot4=12.
\]
Поскольку \(10\ne12\), \key{это не пропорция}.

\subsection{Ещё один способ проверить четыре числа}
Можно ли составить пропорцию \(3{,}06:0{,}9=4{,}08:1{,}2\)?
\[
3{,}06\cdot1{,}2=3{,}672,\qquad
0{,}9\cdot4{,}08=3{,}672.
\]
Да, эта запись является верной пропорцией.

\clearpage
\thispagestyle{empty}
\begin{center}
{\Large\color{darkaccent} Алгоритм: проверка пропорции}
\end{center}
\vspace{16mm}
\begin{center}
\begin{tikzpicture}[node distance=14mm]
\node[flowoval] (start) {Дано равенство \(a:b=c:d\)\ при \(b,d\ne0\)};
\node[flowstep,below=of start] (prod) {Вычислите \(ad\) и \(bc\)};
\node[choice,below=17mm of prod] (judge) {\(ad=bc\)?};
\node[answer,below left=23mm and 19mm of judge] (yes) {Да: равенство\\является пропорцией};
\node[answer,below right=23mm and 19mm of judge] (no) {Нет: равенство\\не является пропорцией};
\draw[arr] (start) -- (prod);
\draw[arr] (prod) -- (judge);
\draw[arr] (judge.south west) -- node[above left,font=\small]{да} (yes.north);
\draw[arr] (judge.south east) -- node[above right,font=\small]{нет} (no.north);
\end{tikzpicture}
\end{center}
\clearpage

\section{Как найти неизвестный член пропорции}

\key{Маршрут:} запишите отношения дробями, перемножьте крайние и средние члены, затем решите полученное уравнение. Не забывайте об ограничениях на знаменатели.

\subsection{Пример 6. Неизвестный средний член}
\[
\frac{0{,}6}{x}=\frac{4}{17},\qquad x\ne0.
\]
Умножим крайние и средние члены:
\[
0{,}6\cdot17=4x,\qquad 10{,}2=4x,
\]
\[
\boxed{x=10{,}2:4=2{,}55.}
\]
Проверка: \(0{,}6:2{,}55=\frac{4}{17}\).

\subsection{Пример 7. Знаменатель со скобками}
\[
\frac{1{,}5}{a+0{,}3}=\frac{6{,}3}{0{,}21},\qquad a\ne-0{,}3.
\]
Весь знаменатель \(a+0{,}3\) --- \key{одно выражение}, поэтому скобки обязательны:
\begin{align*}
1{,}5\cdot0{,}21&=6{,}3(a+0{,}3),\\
0{,}315&=6{,}3(a+0{,}3),\\
a+0{,}3&=0{,}315:6{,}3=0{,}05,\\
a&=\boxed{-0{,}25}.
\end{align*}
Проверка: \(a+0{,}3=0{,}05\ne0\), а обе дроби равны \(30\).

\section{Проценты через пропорцию}

\key{Основная мысль:} целому соответствует \(100\%\), а части --- искомое количество процентов. Записывайте числа под числами, а проценты под процентами.

\subsection{Пример 8. Продолжительность дня}
В сутках \(24\) часа. Продолжительность светового дня --- \(8{,}4\) часа. Сколько процентов суток это составляет?
\[
\begin{array}{rcl}
24\text{ ч} &\longleftrightarrow&100\%,\\
8{,}4\text{ ч} &\longleftrightarrow&x\%.
\end{array}
\]
Составим пропорцию:
\[
\frac{8{,}4}{24}=\frac{x}{100}.
\]
Умножим крайние и средние члены:
\[
24x=8{,}4\cdot100=840,\qquad
\boxed{x=35.}
\]
\key{Ответ:} \(35\%\) суток. Проверка: \(0{,}35\cdot24=8{,}4\) часа.

\key{Общая формула:} если часть \(A\) входит в целое \(B>0\), то её доля в процентах равна
\[
\boxed{p=\frac{A}{B}\cdot100\%.}
\]
Но для текстовых задач особенно полезно понимать именно \key{составление пропорции}: \(B\leftrightarrow100\%\), \(A\leftrightarrow p\%\).

\section{Типичные ошибки и самопроверка}
\begin{itemize}
\item \key{Перепутать доли и числа:} в отношении \(2:7\) числа \(2\) и \(7\) задают количества долей, а не сами части.
\item \key{Не сложить коэффициенты:} одну долю находят делением на сумму коэффициентов отношения.
\item \key{Неверно выразить неизвестное:} из \(4x=10{,}2\) следует \(x=10{,}2:4\), а не \(10{,}2\cdot4\).
\item \key{Потерять скобки:} в дроби \(\frac{1{,}5}{a+0{,}3}\) знаменатель целиком равен \(a+0{,}3\).
\item \key{Пропустить ограничение:} знаменатель в пропорции не может быть нулём.
\item \key{Ожидать только целые доли:} одна доля может быть дробной. Например, \(693:5=138{,}6\).
\item \key{Перепутать часть и целое:} в задаче о процентах всегда сначала установите, чему соответствует \(100\%\).
\end{itemize}

\clearpage
\section*{Тренировочный блок}
\textcolor{muted}{Путь А · 10 задач · от базового уровня к повышенному.}\par
\textit{Решайте последовательно. В задачах с дробями учитывайте, что знаменатель не равен нулю.}
\begin{enumerate}
\item Разделите число \(48\) в отношении \(1:3\).
\item Разделите число \(77\) в отношении \(4:7\).
\item Разделите число \(42\) в отношении \(1:2:4\).
\item Разделите число \(65\) в отношении \(2:3\). Проверьте сумму частей.
\item Докажите, что \(\dfrac35=\dfrac{12}{20}\) является пропорцией.
\item Проверьте, является ли пропорцией равенство \(\dfrac47=\dfrac{18}{28}\). Обоснуйте вывод.
\item Найдите \(x\) из пропорции \(\dfrac5x=\dfrac{15}{21}\).
\item Найдите \(x\) из пропорции \(\dfrac{0{,}8}{x+0{,}5}=\dfrac4{15}\).
\item Продолжительность светового дня --- \(7{,}2\) часа. Какой процент суток она составляет? Решите через пропорцию.
\item Три числа относятся как \(2:5:8\). Самое большое число на \(36\) больше самого маленького. Найдите все три числа и их сумму.
\end{enumerate}
\vfill
\begin{center}
\small\textcolor{muted}{Проверяйте каждый результат: по сумме частей, равенству произведений или подстановкой.}
\end{center}

\clearpage
\section*{Ответы к тренировочному блоку}
\begin{center}
\renewcommand{\arraystretch}{1.55}
\begin{tabularx}{\linewidth}{@{}p{15mm}X@{}}
\toprule
\textbf{№}&\textbf{Ответ}\tabularnewline
\midrule
1 & \(12\) и \(36\).\\
2 & \(28\) и \(49\).\\
3 & \(6,\ 12,\ 24\).\\
4 & \(26\) и \(39\); сумма \(65\).\\
5 & Да: \(3\cdot20=5\cdot12=60\).\\
6 & Нет: \(4\cdot28=112\ne126=7\cdot18\).\\
7 & \(x=7\).\\
8 & \(x=2{,}5\) (\(x\ne-0{,}5\)).\\
9 & \(30\%\).\\
10 & \(12,\ 30,\ 48\); сумма \(90\).\\
\bottomrule
\end{tabularx}
\end{center}
\vspace{12mm}
\section*{Итоговая памятка}
\begin{center}
\renewcommand{\arraystretch}{1.4}
\begin{tabularx}{\linewidth}{@{}p{62mm}X@{}}
\toprule
\textbf{Сигнал в задаче} & \textbf{Полезное действие}\\
\midrule
«Разделить в отношении» & Сложите коэффициенты, найдите одну долю, умножьте на каждый коэффициент.\\
«Верна ли пропорция?» & Сравните произведения крайних и средних членов.\\
«Найдите неизвестный член» & Используйте перекрёстное умножение и решите уравнение.\\
«Сколько процентов?» & Целому поставьте в соответствие \(100\%\); составьте пропорцию.\\
\bottomrule
\end{tabularx}
\end{center}
\vfill
\begin{center}
\small\textcolor{muted}{Лёвин Артём Александрович эксклюзивно для Ксении Васильченко}
\end{center}
\end{document}
